\section{Bounded CNF Expressions}

A \emph{conjunctive normal form} (CNF) expression is any product
$c_1 c_2 \cdots c_r$ of clauses where each clause $c_i$ is a sum
$q_1 + q_2 + \cdots + q_j$ of literals. A \emph{literal} $q$ is
either a variable $p$ or the negation $\bar{p}$ of a variable. For
example, $(p_1 + p_2)(\bar{p}_1 + \bar{p}_2 + p_3)$ is a CNF
expression. For a positive integer $k$ a \emph{$k$-CNF expression} is
a CNF expression where each clause is the sum of at most $k$
literals.

For CNF expressions in general we do not know of any polynomial time
deduction procedure with respect to our learning protocol. The
question of whether one exists is tantalizing because of its apparent
simplicity.

In this section we shall prove that for any fixed integer $k$ the
class of $k$-CNF expressions is learnable. In fact it is learnable
easily in the sense that calls of EXAMPLES suffice and no calls of
ORACLE are required.

In this and subsequent sections we shall use the relation $\Rightarrow$
on concepts as follows: $F \Rightarrow G$ means that for all vectors
$\mathbf{v}$ whenever $F(\mathbf{v}) = 1$ it is also the case that
$G(\mathbf{v}) = 1$. (N.B.\ This is equivalent to the relation
$F \Rightarrow G$ when $F$, $G$ are regarded as functions.) For
brevity we shall often denote both expressions and even vectors by
the concepts they represent. Thus the concept represented by vector
$\mathbf{w}$ is true for exactly those vectors that agree with
$\mathbf{w}$ on all variables on which $\mathbf{w}$ is determined.
The concept represented by expression $f$ is obtained by considering
the Boolean function computed by $f$ and extending its domain to
become a concept. In this section we regard a clause $c_i$ to be a
special kind of expression. In Sections 7 and 8 we consider monomials
$m$, simple products of literals, as special forms of expressions. In
all cases statements of the form $f \Rightarrow g$,
$\mathbf{v} \Rightarrow c_i$, $\mathbf{v} \Rightarrow m$ or
$m \Rightarrow f$ can be understood by interpreting the two sides as
concepts or functions in the obvious way.

\begin{theorem}[Theorem A]
For any positive integer $k$ the class of $k$-CNF expressions is
learnable via an algorithm $A$ that uses
$L = L(h, (2t)^{k+1})$ calls of \textup{EXAMPLES} and no calls of
\textup{ORACLE}, where $t$ is the number of variables.
\end{theorem}

\begin{proof}
The algorithm is initialized with formula $g$ as the product of all
possible clauses of up to $k$ literals from
$\{p_1, \bar{p}_1, p_2, \bar{p}_2, \ldots, p_t, \bar{p}_t\}$.
Clearly the number of ways of choosing a clause of exactly $k$
literals is at most $(2t)^k$ and hence the number of ways of choosing
a clause with up to $k$ literals is less than
$2t + (2t)^2 + \cdots + (2t)^k < (2t)^{k+1}$. This bounds the
initial number of clauses in $g$.

The algorithm then calls EXAMPLES $L$ times to give positive examples
of the concept represented by $k$-CNF formula $f$. For each vector
$\mathbf{v}$ so obtained it deletes all of the remaining clauses in
$g$ that do not contain a literal that is determined to be true in
$\mathbf{v}$. (I.e.\ in the clauses deleted each literal is either
not determined or is negated in $\mathbf{v}$.) More precisely the
following is repeated $L$ times:

\begin{tabbing}
\textbf{begin} \= $\mathbf{v} := \text{EXAMPLES}$ \\
\> for each $c_i$ in $g$ delete $c_i$ if $\mathbf{v} \not\Rightarrow c_i$. \\
\textbf{end}
\end{tabbing}

The claim is that with high probability the value of $g$ after the
$L$-th execution of this block will be the desired approximation to
the formula $f$ that is being learnt.

Initially the set of vectors $\{\mathbf{v} \mid \mathbf{v} \Rightarrow g\}$
is clearly empty. We first claim that as the algorithm proceeds the
set $\{\mathbf{v} \mid \mathbf{v} \Rightarrow g\}$ will always be a
subset of $\{\mathbf{v} \mid \mathbf{v} \Rightarrow f\}$. To prove
this it is clearly sufficient to prove the same statement when both
sets are restricted to only total vectors. Let $B$ be the product of
all the clauses $c$ containing up to $k$ literals with the property
that ``if $\mathbf{v} \Rightarrow f$ then
$\mathbf{v} \Rightarrow c$.'' It is clear that in the course of the
algorithm no clause in $B$ can be ever deleted. Hence it is certainly
true that
$\{\mathbf{v} \mid \mathbf{v} \Rightarrow g\} \subseteq \{\mathbf{v} \mid \mathbf{v} \Rightarrow B\}$.
To establish the claim it remains only to prove that $B$ computes the
same Boolean function as $f$. (In fact $B$ will be the maximal
$k$-CNF formula equivalent to $f$.) It is easy to see that
$f \Rightarrow B$ since by the definition of $B$, for every $c$ in
$B$ it is the case that ``for all $\mathbf{v}$ if
$\mathbf{v} \Rightarrow f$ then $\mathbf{v} \Rightarrow c$.'' To
verify the converse, that $B \Rightarrow f$, we first note that, by
definition, $f$ is a $k$-CNF formula. If some clause in $f$ did not
occur in $B$ then this clause $c'$ would have the property that
``$\exists \mathbf{v}$ such that $\mathbf{v} \Rightarrow f$ but
$\mathbf{v} \not\Rightarrow c'$.'' But this is impossible since if
$c'$ is a clause of $f$ and if $\mathbf{v} \not\Rightarrow c'$ then
$\mathbf{v} \not\Rightarrow f$. We conclude that every $k$-CNF
representation of the function that $f$ represents consists of a set
of clauses that is a subset of the clauses in $B$. Hence
$B \Rightarrow f$.

Let $X = \Sigma D(\mathbf{v})$ with summation over all (not
necessarily total) vectors $\mathbf{v}$ such that
$\mathbf{v} \Rightarrow f$ but $\mathbf{v} \not\Rightarrow g$. This
quantity is defined for each intermediate value of $g$ in the course
of the algorithm and, as is easily verified, it is monotone
decreasing with time. Now clauses will be removed from $g$ whenever
EXAMPLES outputs a vector $\mathbf{v}$ such that
$\mathbf{v} \not\Rightarrow g$. The probability of this happening at
any moment is exactly the current value of $X$. Also, the process of
running the algorithm to completion can have one of just two possible
outcomes: (i) At some point $X$ becomes less than $h^{-1}$, in which
case the $g$ found will approximate to $f$ exactly as required by the
definition of learnability. (ii) The value of $X$ never goes below
$h^{-1}$ and hence $g$ is not an acceptable approximation. The
probability of the latter eventuality is, however, at most $h^{-1}$
since it corresponds to the situation of performing
$L(h, (2t)^{k+1})$ Bernoulli experiments (i.e.\ calls of EXAMPLES)
each with probability greater than $h^{-1}$ of success (i.e.\ finding
a $\mathbf{v}$ such that $\mathbf{v} \not\Rightarrow g$) and
obtaining fewer than $(2t)^{k+1}$ successes (a success being
manifested by the removal of at least one clause from $g$).
\qed
\end{proof}

In conclusion we observe that a CNF expression $g$ immediately yields
a program for computing the associated concept. Given a vector
$\mathbf{v}$ we substitute the determined truth values in $g$. The
concept will be true for $\mathbf{v}$ if and only if all the clauses
are made true by the substitution.
